\section{Attention as Evidence Gathering}
\label{sec:attention-and}

Within the BP interpretation, the attention mechanism performs the
gather phase of message passing: it routes neighboring beliefs into the
residual stream before the update step executes. This corresponds to the
evidence-gathering phase of Pearl's sum-product
algorithm~\cite{pearl1988probabilistic}.

The conjunctive structure is architectural rather than localized inside
any individual attention head. Consider a single transformer layer. Each
attention head scans token positions, concentrates on a relevant
neighbor, and copies information from that neighbor into a designated
region of the current token's residual stream. Only after all heads have
finished does the feed-forward network execute.

The conjunction is therefore not inside any one head. It emerges at the
level of the residual stream, where multiple required evidence sources
become simultaneously available in a shared workspace before the update
step runs. The FFN has no mechanism to execute on partial information:
it reads the entire residual stream state. In this sense, the
architecture enforces a gather-before-update semantics analogous to
conjunctive evidence integration.

\subsection*{The Formal Weight Construction}

The constructive mapping from
\cite{coppola2026transformers} realizes the gather step exactly. Each
attention head uses two sparse weight families:

\begin{itemize}
\item $\mathrm{projectDim}(d)$: a diagonal projection extracting
dimension $d$. Used as $W_Q$ and $W_K$. The resulting $Q \cdot K$ dot
product peaks when token $k$'s stored index matches token $j$'s query
value, implementing exact index matching.

\item $\mathrm{crossProject}(s, d)$: an off-diagonal projection reading
source dimension $s$ and writing to destination dimension $d$. Used as
$W_V$. Routes neighbor $k$'s belief from one residual-stream dimension
into another.
\end{itemize}

At sufficiently high temperature, the softmax concentrates almost
entirely on the matching neighbor. In the constructive setting, sharp
attention behaves like a routing or lookup operation rather than an
inference computation. The probabilistic update occurs in the FFN stage.

\subsection*{A Concrete Example}

Consider three propositions:
\texttt{like(jack,jill)} with belief $m_0 = 0.8$,
\texttt{like(jill,jack)} with belief $m_1 = 0.4$, and
\texttt{date(jack,jill)} whose belief is to be updated.

Head~0 attends to \texttt{like(jack,jill)} and writes $0.8$ into one
region of the \texttt{date} token's residual stream. Head~1 attends to
\texttt{like(jill,jack)} and writes $0.4$ into another region. By the
time the FFN executes, both evidence sources are simultaneously
available. The gather phase is complete before the update step begins.

\subsection*{Evidence from Trained Models}

Mechanistic interpretability studies have identified attention patterns
consistent with the routing behavior predicted by the constructive
mapping. Wang et al.~\cite{wang2022ioi} identify ``name mover heads'' in
GPT-2 small that copy token-specific information to target positions,
behavior closely resembling crossProject-style routing. Their Q/K
circuits also exhibit approximately rank-1 structure, qualitatively
consistent with projectDim-style matching.

Empirical classification studies further identify several recurring head
types:

\begin{itemize}
\item \textbf{Semantic routing heads}: attend to semantically related
tokens determined by content rather than position. These resemble the
gather step of the constructive BP mapping.

\item \textbf{Hub heads}: attend consistently to fixed positions,
especially the BOS token. These appear to distribute or collect global
context information through a shared communication channel.
\end{itemize}

These observations do not by themselves prove the BP interpretation, but
they are structurally consistent with a routing-plus-update view of
transformer computation.

\subsection*{Scaling}

For Llama~3 405B, there are $L \times H = 126 \times 128 = 16{,}128$
distinct attention-head routing patterns. The number of routing
operations per forward pass scales as
$L \times H \times W = 132\mathrm{M}$ at sequence length
$W = 8192$. Within the BP interpretation, these routing operations form
the communication structure of the implicit factor graph.